% Chapter — Optical sensors: star tracker and sun sensor (FR-23). Follows
% the mandatory FR-29 six-section template: purpose; assumptions and domain
% of validity; derivation; discretization and implementation; validation
% evidence; references. Written ahead of the implementation (Phase 6): this
% chapter is the binding specification for the optical sensor modules and
% for the velocity-aberration correction shared with the camera hook
% (Chapter ch:camera). The implementing change adds the chapter-manifest
% entry and the \testid{} evidence table. The equation labels
% eq:optical:aberration, eq:optical:beta, eq:optical:rho, eq:optical:qab,
% eq:optical:stmodel, eq:optical:noiseq, and eq:optical:sunsensor are
% intended to be echoed verbatim in the module source (FR-29 traceability).
\chapter{Sensors: Star Tracker and Sun Sensor}
\label{ch:sensors-optical}

\section{Purpose}
\label{sec:optical:purpose}

This chapter specifies the FR-23 optical attitude sensors: a star tracker
that emits a small-angle-corrupted attitude quaternion with Sun/body
exclusion and slew-rate gating, and a sun sensor that emits a noisy unit
vector toward the apparent Sun. It also defines, once for the whole
project, the \emph{velocity-aberration correction} that FR-23 applies to
every optical truth direction (star tracker, sun sensor, and the camera
hook of Chapter~\ref{ch:camera}): the project's light-time/aberration
policy is that dynamics use geometric ephemeris positions
(Chapter~\ref{ch:ephemeris}), while optical sensor truth adds the
velocity aberration of Section~\ref{sec:optical:aberration} --- a
$\approx \qty{20.5}{\arcsecond}$ effect at Earth's heliocentric orbital
speed, an order of magnitude above an arcsecond-class star-tracker error
budget --- with the residual light-time effects bounded in
Section~\ref{sec:optical:domain}. Phase~6 exit criterion~9 gates an
independent recomputation of the stated aberration formula to
$< \qty{1}{\mas}$, and exit criterion~1 gates the star-tracker error
statistics inside two-sided 95\,\% chi-square bounds over 1{,}000 draws,
so every convention here --- frames, sign, normalization, composition
order --- is stated exactly.

\section{Assumptions and domain of validity}
\label{sec:optical:domain}

Assumptions:
\begin{enumerate}
  \item \textbf{Light-time/aberration policy (FR-23).} Source directions
    are computed from \emph{geometric} (instantaneous) ephemeris
    positions and then corrected for velocity (stellar) aberration only.
    The neglected planetary-aberration term --- the source's own motion
    during the light time --- displaces the apparent Sun by
    $\lesssim \qty{11}{\mas}$ (solar barycentric velocity
    $\lesssim \qty{16}{\metre\per\second}$ acting over
    \qty{499}{\second}) and a Moon-fixed landmark seen from
    \qty{4e8}{\metre} by $\lesssim \qty{0.7}{\arcsecond}$ (lunar orbital
    velocity $\approx \qty{1}{\kilo\metre\per\second}$ over
    \qty{1.3}{\second}); both are below the arcsecond-class sensor noise
    this suite targets, and both bounds scale linearly with source
    velocity and light time for other geometries.
  \item \textbf{First-order aberration.} The first-order formula of
    equation~\eqref{eq:optical:aberration} is used, not the exact
    special-relativistic form; the difference is
    $(\beta^2/4)\sin 2\vartheta + O(\beta^3)$ in apparent angle
    ($\vartheta$ the source--velocity angle), at most
    $\qty{0.52}{\mas}$ at $\beta = 10^{-4}$
    ($v = \qty{30}{\kilo\metre\per\second}$) --- below every stated gate.
  \item \textbf{Stars at infinity.} Catalog directions are truth by
    construction; parallax is not modeled (the simulator has no star
    catalog --- the tracker abstraction below absorbs it).
  \item \textbf{Narrow-field tracker abstraction.} The star tracker is
    modeled at the attitude level: matching, centroiding, and the star
    field itself are abstracted into a body-frame small-angle error with
    configured per-axis standard deviations (about-boresight typically
    the largest). The aberration of the observed star field enters as
    the rigid field rotation of equation~\eqref{eq:optical:rho}; the
    residual differential aberration across the field of view is
    $O(\beta \alpha^2)$ for FOV half-angle $\alpha$ --- below
    \qty{1}{\arcsecond} for $\alpha \le 10^{\circ}$ at
    $\beta = 10^{-4}$ --- and is absorbed in the configured noise.
  \item \textbf{Exclusion angles as configuration.} Exclusion gating
    compares the boresight to the apparent direction of each configured
    excluded body against a configured half-angle per body; the angle is
    understood to include the body's angular radius plus stray-light
    margin. The model does not compute limb geometry for the tracker.
\end{enumerate}

Domain of validity: $\beta = \norm{\vv{v}_{\mathrm{obs}}}/c \le 10^{-3}$
(the first-order error bound above grows as $\beta^2$), configured error
sigmas $\le \qty{e-2}{\radian}$ (the small-angle statistics of
Section~\ref{sec:optical:stats} degrade as $O(\sigma^2)$ beyond that),
and exclusion/FOV half-angles in $(0, \pi/2)$. Out-of-domain
configurations are rejected by the FR-15 validator; the in-loop model is
a total function and never throws. Milliarcsecond-grade astrometry
(relativistic light deflection, parallax, proper motion) is explicitly
out of scope.

\section{Derivation}
\label{sec:optical:derivation}

\subsection{Velocity aberration}
\label{sec:optical:aberration}

Let $\hat{\vv{u}}^{I}$ be the geometric unit direction \emph{from the
observer to the source}, resolved in GCRF axes, and let
\begin{equation}
  \vv{\beta} = \frac{\vv{v}_{\mathrm{obs}}^{I}}{c},
  \qquad
  \vv{v}_{\mathrm{obs}}^{I} = \vv{v}^{I}
    + \vv{v}_{\mathrm{cb}/\mathrm{SSB}}^{I},
  \label{eq:optical:beta}
\end{equation}
where $\vv{v}^{I}$ is the vehicle's truth velocity relative to the
central body (the logged \texttt{truth} channel) and
$\vv{v}_{\mathrm{cb}/\mathrm{SSB}}^{I}$ is the central body's velocity
relative to the solar-system barycenter, assembled from the DE440
segment chain of Chapter~\ref{ch:ephemeris} (for Earth: EMB relative to
the SSB plus Earth relative to the EMB), all in GCRF axes. Aberration is
kinematic --- the observed direction depends on the observer's velocity
in the frame in which the catalog directions are defined, which for
star catalogs (ICRS) is the barycentric frame --- hence the barycentric
composition; using the planet-relative velocity alone would omit the
dominant $\approx \qty{30}{\kilo\metre\per\second}$ annual term.

The apparent direction is, to first order in $\beta$
(Kaplan~\cite{kaplan2005}),
\begin{equation}
  \boxed{\;
  \hat{\vv{u}}'
    = \frac{\hat{\vv{u}} + \vv{\beta}
        - (\hat{\vv{u}} \cdot \vv{\beta})\,\hat{\vv{u}}}
      {\bigl\lVert \hat{\vv{u}} + \vv{\beta}
        - (\hat{\vv{u}} \cdot \vv{\beta})\,\hat{\vv{u}} \bigr\rVert} ,
  \;}
  \label{eq:optical:aberration}
\end{equation}
i.e.\ the geometric direction plus the component of $\vv{\beta}$
perpendicular to it, renormalized. The sign convention, stated
explicitly because it is the usual defect: the \emph{apparent} direction
is displaced from the geometric direction \emph{toward the observer's
velocity vector} (the telescope tilts forward), by the angle
\begin{equation}
  \delta\vartheta = \beta \sin\vartheta + O(\beta^2),
  \qquad
  \vartheta = \angle(\hat{\vv{u}}, \vv{\beta}),
  \label{eq:optical:abmag}
\end{equation}
whose maximum $\beta = v_{\mathrm{obs}}/c$ evaluates to
\qty{20.49}{\arcsecond} at Earth's mean heliocentric speed
$v = \qty{29.78}{\kilo\metre\per\second}$ and
\qty{20.64}{\arcsecond} at \qty{30}{\kilo\metre\per\second} --- the
"$\approx \qty{20.5}{\arcsecond}$" of FR-23. Equation
\eqref{eq:optical:aberration} is \emph{the} normative formula: exit
criterion~9 recomputes it independently (NumPy, from the logged truth
and ephemeris inputs) and gates agreement to $< \qty{1}{\mas}$; the
formula-choice error relative to the exact relativistic form is bounded
in Section~\ref{sec:optical:domain}, item~2, and was verified
numerically over the full $\vartheta$ range during drafting.

\subsection{Star tracker measurement model}
\label{sec:optical:startracker}

The tracker is configured with a body-frame boresight unit vector
$\hat{\vv{b}}^{B}$, per-axis error standard deviations
$\vv{\sigma}_{\mathrm{st}} = (\sigma_1, \sigma_2, \sigma_3)$
(\unit{\radian}, body axes), exclusion half-angles
$\theta_{\mathrm{excl},j}$ for the Sun and each configured excluded
body, and a slew-rate limit $\omega_{\max}$.

\paragraph{Aberration as a field rotation.} Over a narrow field of view
the aberration field of equation~\eqref{eq:optical:aberration} acts, to
first order, as a rigid rotation of the observed sky: for directions
$\hat{\vv{u}}$ near the boresight $\hat{\vv{b}}^{I} =
{\dcm{I}{B}}^{\mathsf T}\hat{\vv{b}}^{B}$,
\begin{equation}
  \hat{\vv{u}}' = \hat{\vv{u}} + \vv{\rho} \times \hat{\vv{u}}
    + O(\beta\,\alpha^2),
  \qquad
  \boxed{\;
  \vv{\rho}^{I} = \hat{\vv{b}}^{I} \times \vv{\beta}^{I},
  \;}
  \label{eq:optical:rho}
\end{equation}
verified at the boresight: $\vv{\rho} \times \hat{\vv{b}} =
\vv{\beta} - (\hat{\vv{b}} \cdot \vv{\beta})\hat{\vv{b}}$, exactly the
first-order displacement of equation~\eqref{eq:optical:aberration}. A
tracker that matches this rotated sky against the geometric catalog
reports the attitude of the body relative to the \emph{apparent}
inertial frame. Working in the project convention
(Chapter~\ref{ch:notation}): the apparent frame $I'$ satisfies
$\dcm{I'}{I} = \mat{I}_3 + \skewm{\vv{\rho}}$, so the reported attitude
is $\quat{q}_{i'2b} = \quat{q}_{i'2i} \quatmul \quat{q}_{i2b}$ with
$\dcm{I'}{I}(\quat{q}_{i'2i})$ as above; by
equation~\eqref{eq:notation:quat2dcm} the corresponding unit quaternion
has vector part $-\vv{\rho}/2$ to first order. The model uses the exact
unit quaternion of the rotation vector $-\vv{\rho}$:
\begin{equation}
  \quat{q}_{\mathrm{ab}}
    = \Bigl[\cos\tfrac{\rho}{2},\;
        -\sin\tfrac{\rho}{2}\,\frac{\vv{\rho}}{\rho}\Bigr],
  \qquad \rho = \norm{\vv{\rho}},
  \qquad
  \quat{q}_{\mathrm{ab}} = [1, \vv{0}] \text{ when } \rho = 0 .
  \label{eq:optical:qab}
\end{equation}
Consistency check, binding on the implementation: the boresight
direction implied by the aberrated attitude,
${\dcm{I}{B}(\quat{q}_{\mathrm{ab}} \quatmul
\quat{q}_{i2b})}^{\mathsf T} \hat{\vv{b}}^{B} = \hat{\vv{b}}^{I} -
\vv{\rho} \times \hat{\vv{b}}^{I} + O(\rho^2)$, points \emph{away} from
the velocity --- correct, because the star whose apparent position sits
on the boresight has its catalog position displaced against the
velocity.

\paragraph{Measurement equation.} With the noise quaternion
$\delta\quat{q}_n$ below, the emitted measurement is
\begin{equation}
  \boxed{\;
  \quat{q}_{\mathrm{meas}}
    = \quat{q}_{\mathrm{ab}} \quatmul \quat{q}_{i2b}^{\mathrm{true}}
      \quatmul \delta\quat{q}_n ,
  \;}
  \label{eq:optical:stmodel}
\end{equation}
Hamilton, scalar-first, $\quat{q}_{i2b}$ chaining per decision D-7 and
equation~\eqref{eq:notation:quatcomp}: the aberration is an
inertial-side factor (left), the sensor noise a body-side factor
(right). Transcription note: in the composition convention of Markley
and Crassidis~\cite{markley2014}, where $\mat{A}(\quat{p} \quatmul
\quat{q}) = \mat{A}(\quat{p})\mat{A}(\quat{q})$, the same physical model
reads with the factor order reversed
($\delta\quat{q} \quatmul \quat{q}_{\mathrm{true}}$ for a body-side
error); the form~\eqref{eq:optical:stmodel} is the normative one for
this project's code, logs, and cross-checks.

\paragraph{Noise quaternion.} Draw
$\vv{\epsilon} \sim \mathcal{N}\bigl(\vv{0},\,
\mathrm{diag}(\sigma_1^2, \sigma_2^2, \sigma_3^2)\bigr)$
(body axes) and form the \emph{exact} rotation
\begin{equation}
  \delta\quat{q}_n
    = \Bigl[\cos\tfrac{\theta}{2},\;
        \sin\tfrac{\theta}{2}\,\frac{\vv{\epsilon}}{\theta}\Bigr],
  \qquad \theta = \norm{\vv{\epsilon}},
  \qquad
  \delta\quat{q}_n = [1, \vv{0}] \text{ when } \theta = 0 ,
  \label{eq:optical:noiseq}
\end{equation}
the exponential map of the drawn rotation vector. Using the exact map
rather than a small-angle construction costs nothing and makes the
statistic of Section~\ref{sec:optical:stats} exactly chi-square
distributed, because the extraction below recovers $\vv{\epsilon}$
identically.

\paragraph{Validity gating.} The emitted flag is
\begin{equation}
  \mathrm{valid} =
    \Bigl[\bigwedge_{j}
      \angle\bigl(\hat{\vv{b}}^{I},\, \hat{\vv{u}}'_j\bigr)
        \ge \theta_{\mathrm{excl},j}\Bigr]
    \wedge
    \bigl[\norm{\vv{\omega}^{B}} \le \omega_{\max}\bigr],
  \label{eq:optical:gating}
\end{equation}
where $j$ ranges over the Sun and the configured excluded bodies,
$\hat{\vv{u}}'_j$ is the \emph{apparent} (aberrated) direction to body
$j$'s center per equation~\eqref{eq:optical:aberration} (the
$\le \qty{20.5}{\arcsecond}$ geometric-versus-apparent difference is
immaterial against degree-scale exclusion angles, but a single code path
computes all optical directions), $\angle(\vv{a},\vv{b}) =
\operatorname{atan2}(\norm{\vv{a}\times\vv{b}}, \vv{a}\cdot\vv{b})$, and
$\vv{\omega}^{B}$ is the true body rate. Semantics, normative: the
measurement is computed, logged, and flagged every sample regardless of
validity; noise draws are unconditional (stream-schedule stability, as
in Chapter~\ref{ch:sensors-imu}); consumers must treat
$\mathrm{valid} = \mathrm{false}$ samples as absent --- the reference
navigation components of Chapters~\ref{ch:gnc-builtin} and~\ref{ch:ekf}
skip the corresponding updates.

\subsection{Sun sensor measurement model}
\label{sec:optical:sunsensor}

The sun sensor is configured with a boresight $\hat{\vv{s}}^{B}$, a
field-of-view half-angle $\alpha_{\mathrm{ss}}$, and an error standard
deviation $\sigma_{\mathrm{ss}}$ (\unit{\radian}). With
$\hat{\vv{u}}'^{\,I}_{\odot}$ the apparent Sun direction from
equation~\eqref{eq:optical:aberration} and $\hat{\vv{u}}^{B} =
\dcm{I}{B}(\quat{q}_{i2b}^{\mathrm{true}})\, \hat{\vv{u}}'^{\,I}_{\odot}$
its body-frame resolution, the measurement is
\begin{equation}
  \boxed{\;
  \hat{\vv{u}}_{\mathrm{meas}}
    = \frac{\hat{\vv{u}}^{B} + \vv{\eta}}
           {\norm{\hat{\vv{u}}^{B} + \vv{\eta}}},
  \qquad
  \vv{\eta} \sim \mathcal{N}\bigl(\vv{0},\,
    \sigma_{\mathrm{ss}}^2\,\mat{I}_3\bigr),
  \;}
  \label{eq:optical:sunsensor}
\end{equation}
the simplest standard unit-vector error model: the radial noise
component drops on normalization and the two tangential components give
a direction error of standard deviation $\sigma_{\mathrm{ss}}$ per axis
to $O(\sigma_{\mathrm{ss}}^2)$. Validity:
\begin{equation}
  \mathrm{valid} =
    \bigl[\angle(\hat{\vv{u}}^{B}, \hat{\vv{s}}^{B})
      \le \alpha_{\mathrm{ss}}\bigr]
    \wedge \bigl[\nu > 0\bigr],
  \label{eq:optical:sungate}
\end{equation}
with $\nu$ the illumination fraction of the Chapter~\ref{ch:srp} conical
shadow model (equation~\eqref{eq:srp:nu}): a sensor in total umbra sees
no Sun; partial illumination ($0 < \nu < 1$) is treated as visible, with
the centroid shift of a partially occulted disk left unmodeled (bounded
by the solar angular radius, and transited in seconds during eclipse
entry/exit).

\subsection{Acceptance statistics}
\label{sec:optical:stats}

Exit criterion~1 gates the star tracker over $M = 1000$ seeded draws at
fixed truth. Per draw $i$, the evaluation removes the deterministic
factors of equation~\eqref{eq:optical:stmodel} and extracts the error
rotation vector exactly:
\begin{equation}
  \delta\quat{q}^{(i)}
    = \bigl(\quat{q}_{\mathrm{ab}} \quatmul
        \quat{q}_{i2b}^{\mathrm{true}}\bigr)^{-1}
      \quatmul \quat{q}_{\mathrm{meas}}^{(i)},
  \qquad
  \hat{\vv{\epsilon}}^{(i)}
    = \hat\theta\,\hat{\vv{u}},\quad
  \hat\theta = 2\operatorname{atan2}\bigl(\norm{\delta\vv{q}_v},
    |\delta q_w|\bigr),\quad
  \hat{\vv{u}} = \operatorname{sgn}(\delta q_w)\,
    \frac{\delta\vv{q}_v}{\norm{\delta\vv{q}_v}},
  \label{eq:optical:extract}
\end{equation}
with $\operatorname{sgn}(0) = +1$ and $\hat{\vv{\epsilon}} = \vv{0}$
when $\norm{\delta\vv{q}_v} = 0$. By construction
$\hat{\vv{\epsilon}}^{(i)} = \vv{\epsilon}^{(i)}$ (the log map inverts
equation~\eqref{eq:optical:noiseq}; the $\theta > \pi$ wrap has
probability zero for in-domain sigmas), so
\begin{equation}
  q^{(i)} = \hat{\vv{\epsilon}}^{(i)\mathsf T}\,
    \mathrm{diag}(\sigma_1^2, \sigma_2^2, \sigma_3^2)^{-1}\,
    \hat{\vv{\epsilon}}^{(i)} \sim \chi^2_3
  \quad\text{exactly, i.i.d.},
  \label{eq:optical:ststat}
\end{equation}
and the gate is two-sided on the ensemble mean:
\begin{equation}
  \frac{\chi^2_{0.025}(3M)}{M}
    \le \frac{1}{M}\sum_{i=1}^{M} q^{(i)}
    \le \frac{\chi^2_{0.975}(3M)}{M},
  \qquad
  [2.850,\; 3.154] \text{ for } M = 1000 ,
  \label{eq:optical:stbounds}
\end{equation}
quantiles per the Wilson--Hilferty evaluation of
Chapter~\ref{ch:ekf}. The sun sensor uses the tangent-plane statistic:
with the deterministic orthonormal basis $\hat{\vv{e}}_1 =
(\hat{\vv{u}}^{B} \times \hat{\vv{z}})/\norm{\cdot}$ (falling back to
$\hat{\vv{x}}$ when $\norm{\hat{\vv{u}}^{B} \times \hat{\vv{z}}} <
10^{-8}$), $\hat{\vv{e}}_2 = \hat{\vv{u}}^{B} \times \hat{\vv{e}}_1$,
the components $e_a^{(i)} = \hat{\vv{e}}_a \cdot
\hat{\vv{u}}_{\mathrm{meas}}^{(i)}$ give $q^{(i)} = \bigl((e_1^2 +
e_2^2)/\sigma_{\mathrm{ss}}^2\bigr)^{(i)} \sim \chi^2_2$ to
$O(\sigma_{\mathrm{ss}}^2)$, gated on the mean within
$[\chi^2_{0.025}(2M), \chi^2_{0.975}(2M)]/M = [1.878,\; 2.126]$ for
$M = 1000$.

\section{Discretization and implementation}
\label{sec:optical:implementation}

Binding implementation notes for the Phase~6 modules:
\begin{enumerate}
  \item \textbf{Modules and traceability.} Star tracker and sun sensor
    are core-side \texttt{ISensor} modules sampled on the major-cycle
    schedule (D-5), on the named streams \texttt{sensors.st} and
    \texttt{sensors.sun} respectively (D-9). Source comments echo the
    labels \texttt{eq:optical:aberration}, \texttt{eq:optical:rho},
    \texttt{eq:optical:qab}, \texttt{eq:optical:stmodel},
    \texttt{eq:optical:noiseq}, and \texttt{eq:optical:sunsensor}
    verbatim; the aberration of
    equation~\eqref{eq:optical:aberration} is one shared core function
    used by both sensors and by the camera hook
    (Chapter~\ref{ch:camera}) --- a single code path is what makes exit
    criterion~9 a single gate.
  \item \textbf{Draw schedule, normative.} Star tracker: three normals
    per sample ($\epsilon_1, \epsilon_2, \epsilon_3$). Sun sensor: three
    normals per sample ($\eta_1, \eta_2, \eta_3$). Unconditional draws,
    per the Chapter~\ref{ch:sensors-imu} rule.
  \item \textbf{Quaternion output.} Scalar-first per
    Table~\ref{tab:notation:eigenmap}; the emitted quaternion is
    normalized after the product~\eqref{eq:optical:stmodel} (the factors
    are unit to rounding; the explicit normalization pins the invariant)
    and is \emph{not} sign-canonicalized --- consumers own the double
    cover, as does the extraction~\eqref{eq:optical:extract}.
  \item \textbf{Ephemeris inputs.} $\vv{\beta}$ per
    equation~\eqref{eq:optical:beta} from the same DE440 chain the
    dynamics use; for a Mars- or Moon-centered run the central-body
    barycentric velocity replaces the Earth chain. No sensor code reads
    the clock or any state outside \texttt{GncInput}-visible truth
    supplied by the core scheduler.
  \item \textbf{Logged channels.} Measurement, validity flag, and the
    per-sample apparent-direction inputs land in
    \texttt{sensors.st.*} / \texttt{sensors.sun.*} with unit-suffixed
    names, at the sensor rate.
\end{enumerate}

\section{Validation evidence}
\label{sec:optical:validation}

The modules implementing this chapter are
\texttt{cpp/src/sensors/optical.cpp} (both sensors and the shared
aberration function).

\begin{table}[htbp]
  \centering
  \caption{Validation evidence for the optical sensors (doctest,
    \texttt{cpp/tests/test\_sensors.cpp}; pytest,
    \texttt{tests/python/test\_p6\_optical\_gates.py} and
    \texttt{test\_gnc\_missions.py}).}
  \label{tab:optical:validation}
  \begin{tabular}{lp{8.2cm}}
    \toprule
    Test ID & Acceptance basis \\
    \midrule
    \testid{sensors_aberration_magnitude_and_direction} &
      Equation~\eqref{eq:optical:aberration} against the closed-form
      magnitude~\eqref{eq:optical:abmag}: the perpendicular-source
      displacement is \qty{20.49}{\arcsecond} at
      $v = \qty{29.78}{\kilo\metre\per\second}$ to $10^{-3}$ relative, the
      output is a unit vector, and the displacement is verified to be
      \emph{toward} the velocity (the sign convention this chapter calls
      out as the usual defect). Parallel and antiparallel sources are
      undisplaced; the $\beta\sin\vartheta$ law holds across a nine-point
      $\vartheta$ grid spanning $[0, \pi]$ including both extremes; and
      $\vv{\beta} = \vv{0}$ returns the input bit-unchanged. The
      composition of equation~\eqref{eq:optical:beta} is checked to sum
      the vehicle and central-body barycentric velocities before dividing
      by $c$. \\
    \testid{sensors_startracker_noise_extraction_is_the_drawn_rotation} &
      With the observer at rest (so $\vv{\rho} = \vv{0}$ and
      $\quat{q}_{\mathrm{ab}}$ is the identity) about a non-identity truth
      attitude, the extraction~\eqref{eq:optical:extract} returns the drawn
      rotation vector of equation~\eqref{eq:optical:noiseq} to $10^{-11}$
      relative over eight seeded draws --- the exactness that makes the
      acceptance statistic chi-square rather than approximately so. \\
    \testid{sensors_startracker_chi_square_over_1000_draws} &
      Exit criterion~1: the ensemble-mean
      statistic~\eqref{eq:optical:ststat} over $M = 1000$ seeded draws
      falls inside the~\eqref{eq:optical:stbounds} interval
      $[2.850, 3.154]$, and each of the three per-axis sample variances
      matches its configured sigma within the corresponding two-sided
      95\,\% interval $[0.914, 1.090]$. \\
    \testid{sensors_sunsensor_chi_square_over_1000_draws} &
      The Section~\ref{sec:optical:stats} tangent-plane statistic over
      $M = 1000$ seeded draws falls inside $[1.878, 2.126]$, the two-sided
      95\,\% interval for $\chi^2_2$. \\
    \testid{sensors_factory_kind_vocabulary} &
      Both optical kinds construct from the factory at their configured
      rate and report their own \texttt{sensors.*} group names, so a
      sensor and its log declaration cannot drift apart. \\
    \testid{test_full_sensor_suite_reruns_bit_identical} &
      Exit criterion~1 bit identity, optical half. Two seeded runs of the
      reference mission with all six FR-23 sensor kinds declared --- the
      \texttt{sensors.startracker} and \texttt{sensors.sunsensor}
      declarations among them, both carrying nonzero sigmas, so the noise
      draws of equation~\eqref{eq:optical:noiseq} are live --- agree on the
      whole-file SHA-256 and then channel by channel under array bit
      equality, so each optical stream is pinned individually. \\
    \testid{test_aberration_matches_independent_reference} &
      Exit criterion~9. The logged apparent Sun directions match
      \texttt{tests/refs/aberration.py} --- a NumPy transcription of
      equations~\eqref{eq:optical:beta}--\eqref{eq:optical:abmag} written
      from this chapter alone, importing nothing from the core and
      deriving its own constants --- across 300 records of a committed
      ephemeris-backed scenario, and rotated into body axes by the equally
      independent \texttt{tests/refs/quaternions.quat\_to\_dcm} rather than
      by the core's own DCM. Measured worst residual
      \qty{4.726e-8}{\mas} against a committed assertion of
      \qty{e-5}{\mas}, which implies the criterion's \qty{1}{\mas} by five
      orders of magnitude; see the tolerance note below. \\
    \testid{test_aberration_fixture_can_see_an_attitude_convention_error} &
      Guards the row above against a fixture that cannot detect what it
      substitutes. An angular separation is invariant under a rotation
      applied to both arguments, and at $q_w = 0$ the frame-transformation
      DCM of equation~\eqref{eq:notation:quat2dcm} is exactly symmetric ---
      $\mat{C} - \mat{C}^{T} = -4 q_w \skewm{\vv{q}_v}$ vanishes
      identically --- so on an attitude confined to that plane a transposed
      convention is invisible to the residual no matter which
      implementation supplies the DCM. This asserts the fixture carries
      $\max \lvert \mat{C} - \mat{C}^{T} \rvert > 0.05$; the committed
      off-axis slew reaches \num{0.18}. \\
    \testid{test_aberration_signal_dominates_the_gate} &
      Non-vacuity of the row above: the aberration displacement the gate
      resolves is really present, between \num{20000} and
      \qty{21000}{\mas} on this geometry, so agreement at the gate is a
      statement about the correction rather than about its absence. \\
    \testid{test_star_tracker_statistic_passes_the_reference_gate} &
      The criterion-1 star-tracker statistic re-gated end to end through
      \texttt{tests/refs/sensor\_stats.py}, which runs the aberration
      factor~\eqref{eq:optical:qab} backwards out of the logged quaternion.
      A wrong $\quat{q}_{\mathrm{ab}}$ appears here as a mean offset rather
      than as noise --- coverage the C++ gate cannot supply, since it fixes
      $\vv{v} = \vv{0}$ and so composes against the identity. \\
    \bottomrule
  \end{tabular}
\end{table}

\paragraph{What the criterion-9 tolerance is, and what it is not.} The
\qty{1}{\mas} figure is the \emph{requirement} exit criterion~9 states, not
a measure of what the implementation achieves and not a physical error
budget. Both sides of the comparison evaluate the same normative
equation~\eqref{eq:optical:aberration}, so the residual is an arithmetic
disagreement between two transcriptions, and the measured worst value is
\qty{4.726e-8}{\mas} --- some $2.1\times10^{7}$ times inside the
criterion's own figure.
Setting the assertion at the requirement rather than near the achievable
agreement had a consequence worth stating plainly: a reference mutation
that drops the transverse projection from
equation~\eqref{eq:optical:aberration} (writing $\vv{u} + \vv{\beta}$ for
$\vv{u} + \vv{\beta} - (\vv{u}\cdot\vv{\beta})\vv{u}$, a change at second
order only) moves the residual to \qty{0.4696}{\mas} on this fixture and
therefore passed a \qty{1}{\mas} gate. It passed because that tolerance was
loose relative to the achievable agreement, not because \qty{1}{\mas} is
forced by anything physical. The assertion has accordingly been tightened,
which is the remediation recorded under \texttt{KNOWN-ISSUE-P6-2}: the
committed constant now reads
\texttt{ABERRATION\_TOL\_MAS} \texttt{= 1e-5}, retaining $211.6\times$
headroom over the measured residual while rejecting the drop-transverse
mutation by a factor of about $4.7\times10^{4}$. The criterion's own
\qty{1}{\mas} figure is asserted alongside it, so the suite states the
requirement it meets and not only the tighter bound it is held to.

Tightening the tolerance is necessary but not sufficient, because the
residual is an angular separation and so is invariant under a rotation
applied to both of its arguments. Two further changes make the gate
sensitive to the attitude convention. The reference side now rotates
through \texttt{tests/refs/quaternions.quat\_to\_dcm}: the core's own DCM
previously appeared on both sides of the separation and cancelled exactly,
so no convention error could reach the residual. The fixture's commanded
attitude was also moved off the body $+Z$ axis, because the earlier slew
held $q_w = 0$ for the whole run, and there
$\mat{C} - \mat{C}^{T} = -4 q_w \skewm{\vv{q}_v}$ vanishes identically ---
a transposed convention was invisible by geometry regardless of which
implementation supplied the DCM. The off-axis slew carries
$\lvert q_w \rvert$ up to \num{0.060} and
$\max \lvert \mat{C} - \mat{C}^{T} \rvert$ up to \num{0.18};
\testid{test_aberration_fixture_can_see_an_attitude_convention_error} pins
that asymmetry so the fixture cannot drift back into the blind plane.

Distinct from the tolerance, and genuinely a modelling choice rather than a
gate weakness: this chapter declares the first-order
form~\eqref{eq:optical:aberration} normative, and the gap to the exact
relativistic form is \qty{0.51}{\mas} at Earth's mean heliocentric speed
and \qty{0.81}{\mas} once low-Earth-orbit speed is added. Criterion~9
recomputes the first-order form on both sides deliberately; gating the core
against the exact form would report a deliberate modelling decision as an
implementation error. That gap is measured, and asserted against
assumption~2's own bound, by
\testid{test_first_order_versus_exact_gap_is_recorded}, which is recorded
as non-normative for exactly this reason.

Outstanding, recorded rather than tabled:
\begin{itemize}
  \item \textbf{Attitude convention inside the criterion-9 gate.} The
    reference side of \testid{test_aberration_matches_independent_reference}
    rotates its apparent direction into the body frame with the core's own
    \texttt{quat\_to\_dcm}, and the residual is an angular separation, which
    is invariant under a rotation common to both sides. The core DCM
    therefore cancels exactly and no error in it could be detected there.
    An independently validated implementation exists at
    \texttt{tests/refs/quaternions.py} and is already used by the
    Chapter~\ref{ch:camera} pixel gate; substituting it closes the hole.
    The convention is not unpinned in the project ---
    \testid{test_rotation_bindings_match_goldens} gates that same binding
    against the Chapter~\ref{ch:rotations} goldens to $10^{-15}$ per
    element --- but it is not pinned \emph{here}, and the criterion-9 row
    should not be read as evidence for it.
  \item \textbf{Gating geometry.} Deterministic sweeps of the boresight
    across each exclusion boundary and of the rate across
    $\omega_{\max}$, and an umbra transit dropping the sun-sensor flag,
    each confirming the flag flips at the configured threshold with draws
    consumed identically on both sides. The gating logic of
    equations~\eqref{eq:optical:gating} and~\eqref{eq:optical:sungate} is
    implemented and its draws are unconditional by construction, but the
    boundary sweeps are not yet committed as tests.
\end{itemize}

A scope note on equation~\eqref{eq:optical:gating}: the implemented
exclusion set is the Sun and the central body, each with its own
configured half-angle. The general "each configured excluded body" form
requires a configured third-body table that the sensor layer does not yet
receive; the single-code-path treatment of apparent directions is
unaffected, since both implemented exclusions run through
equation~\eqref{eq:optical:aberration}.

\section{References}
\label{sec:optical:references}

The aberration treatment follows the IAU-resolution exposition of
Kaplan~\cite{kaplan2005} (velocity aberration and its first-order form),
restricted per the FR-23 policy to the stellar-aberration term; the
attitude-measurement abstraction and quaternion conventions follow
Markley and Crassidis~\cite{markley2014}, transcribed to the project
convention of Chapter~\ref{ch:notation} as noted at
equation~\eqref{eq:optical:stmodel}. Full bibliographic details appear
in the bibliography.
